\section{Canonical LR(1) dan LALR(1)}
\label{sec:ch12-lr1}

\subsection{Item LR(1)}

Item LR(1) menambahkan lookahead:
\[
[A \to \alpha \cdot \beta,\ a],
\]
dengan \(a\) terminal atau \(\$\). Lookahead menunjukkan token yang
boleh muncul setelah reduksi. Informasi ini membuat LR(1) lebih
presisi daripada SLR(1), yang memakai FOLLOW secara global.

\subsection{Closure LR(1)}

Jika item \([A \to \alpha \cdot B \beta,\ a]\) berada dalam
himpunan, maka untuk setiap aturan \(B \to \gamma\) dan setiap
\(b \in \text{FIRST}(\beta a)\), tambahkan
\([B \to \cdot \gamma,\ b]\). Perhatikan bahwa lookahead dihitung
dari \(\beta a\), bukan hanya dari FOLLOW(\(B\)).

\subsection{Canonical LR(1)}

Koleksi kanonik LR(1) dibangun seperti LR(0), tetapi state memuat
item dengan lookahead. Jumlah state LR(1) biasanya jauh lebih besar
daripada LR(0). Untuk grammar bahasa nyata, ledakan state dapat
menjadi masalah.

\subsection{LALR(1)}

LALR(1) menggabungkan state LR(1) yang memiliki inti (\eng{core})
sama, yaitu himpunan item tanpa lookahead. Penggabungan ini
mengurangi jumlah state mendekati LR(0), tetapi mempertahankan
sebagian presisi LR(1).

\begin{figure}[H]
\centering
\begin{tikzpicture}[node distance=6mm and 8mm]
  \node[phase, text width=4.2cm] (lr1a) {LR(1): \([A \to \alpha \cdot, a]\)};
  \node[phase, text width=4.2cm, below=of lr1a] (lr1b)
        {LR(1): \([A \to \alpha \cdot, b]\)};
  \node[phase, text width=4.2cm, right=of lr1a] (lalr)
        {LALR: \([A \to \alpha \cdot, \{a,b\}]\)};
  \draw[arrow] (lr1a) -- (lalr);
  \draw[arrow] (lr1b) -- (lalr);
\end{tikzpicture}
\caption{Penggabungan state LR(1) dengan inti sama menjadi LALR(1)}
\label{fig:ch12-lalr}
\end{figure}

\subsection{Perbandingan}

\begin{table}[H]
\centering
\caption{Perbandingan SLR, LR(1), dan LALR}
\label{tab:ch12-perbandingan}
\begin{tabularx}{\textwidth}{@{}lXXX@{}}
\toprule
\textbf{Aspek} & \textbf{SLR(1)} & \textbf{LR(1)} & \textbf{LALR(1)} \\
\midrule
Lookahead & FOLLOW global & Per item & Per inti state \\
Jumlah state & Kecil & Besar & Sedang \\
Kekuatan & Paling lemah & Paling kuat & Hampir LR(1) \\
Konflik & Lebih mungkin & Paling sedikit & Sedikit \\
Praktik & Jarang & Jarang & Umum (Bison/Yacc) \\
\bottomrule
\end{tabularx}
\end{table}

\subsection{Artefak: Tabel ACTION/GOTO LALR}

Langkah membangun tabel LALR(1) adalah:
\begin{enumerate}[label=\arabic*.,leftmargin=*]
  \item bangun koleksi kanonik LR(1) dengan lookahead per item;
  \item gabungkan state yang memiliki inti (\eng{core}) sama;
  \item gabungkan lookahead item yang bersesuaian;
  \item isi ACTION dan GOTO dari item gabungan;
  \item periksa konflik shift-reduce dan reduce-reduce.
\end{enumerate}

Tabel \tabref{tab:ch12-artefak-lalr} menunjukkan format pencatatan
state LALR yang digabung. Untuk grammar
\(S \to A,\ A \to aA \mid b\), penggabungan inti tidak menimbulkan
konflik baru, sehingga tabel LALR sama dengan tabel SLR pada
\tabref{tab:ch12-artefak-slr}.

\begin{table}[H]
\centering
\caption{Format pencatatan state LALR hasil penggabungan inti}
\label{tab:ch12-artefak-lalr}
\begin{tabularx}{\textwidth}{@{}l l l X@{}}
\toprule
\textbf{State LALR} & \textbf{Inti item} & \textbf{Lookahead} &
\textbf{Aksi} \\
\midrule
\(I_2\) & \(\{S \to A \cdot\}\) & \(\{\$\}\) &
reduce \(S \to A\) \\
\(I_4\) & \(\{A \to b \cdot\}\) & \(\{\$\}\) &
reduce \(A \to b\) \\
\(I_5\) & \(\{A \to aA \cdot\}\) & \(\{\$\}\) &
reduce \(A \to aA\) \\
\bottomrule
\end{tabularx}
\end{table}

Bison dan Yacc secara default membangun tabel LALR(1). Jika konflik
masih muncul, opsi LR(1) kanonik dapat digunakan, tetapi dengan
biaya ukuran tabel yang lebih besar.
